%\title{Amrita Vishwa Vidyapeetham Thesis Template}
\RequirePackage[l2tabu]{nag}		% Warns for incorrect (obsolete) LaTeX usage
%
%
% File: memoirthesis.tex
% Author: Gayathri Manikutty
% Description: Contains the thesis template using memoir class,
% which is mainly based on book class but permits better control of 
% chapter styles for example. This template is an adaptation and 
% modification of Oscar's.
% 
% Memoir is a flexible class for typesetting poetry, fiction, 
% non-fiction and mathematical works as books, reports, articles or
% manuscripts. CTAN repository is found at:
% http://www.ctan.org/tex-archive/macros/latex/contrib/memoir/
%
%
%
% AVV guidelines:
%
% The dissertation must be printed on A4 white paper. Paper up to A3 may be used
% for maps, plans, diagrams and illustrative material. Pages (apart from the
% preliminary pages) should normally be double-sided.
%
% Memoir class loads useful packages by default (see manual).
\documentclass[a4paper,11pt,leqno,openbib]{memoir} %add 'draft' to turn draft option on (see below)
%
%
% Adding metadata:
\usepackage{datetime}
\usepackage{ifpdf}
\ifpdf
\pdfinfo{
   /Author (Author's name)
   /Title (PhD Thesis)
   /Keywords (One; Two;Three)
   /CreationDate (D:\pdfdate)
}
\fi
% When draft option is on. 
\ifdraftdoc 
	\usepackage{draftwatermark}				%Sets watermarks up.
	\SetWatermarkScale{0.3}
	\SetWatermarkText{\bf Draft: \today}
\fi
%
% Declare figure/table as a subfloat.
\newsubfloat{figure}
\newsubfloat{table}
% Better page layout for A4 paper, see memoir manual.
\settrimmedsize{297mm}{210mm}{*}
\setlength{\trimtop}{0pt} 
\setlength{\trimedge}{\stockwidth} 
\addtolength{\trimedge}{-\paperwidth} 
\settypeblocksize{634pt}{448.13pt}{*} 
\setulmargins{4cm}{*}{*} 
\setlrmargins{*}{*}{1.5} 
\setmarginnotes{17pt}{51pt}{\onelineskip} 
\setheadfoot{\onelineskip}{2\onelineskip} 
\setheaderspaces{*}{2\onelineskip}{*} 
\checkandfixthelayout
%
\frenchspacing
% Font with math support: New Century Schoolbook
\usepackage{fouriernc}
\usepackage[T1]{fontenc}
%
% AVV guidelines:
%
% Text should be in double or 1.5 line spacing, and font size should be
% chosen to ensure clarity and legibility for the main text and for any
% quotations and footnotes. Margins should allow for eventual hard binding.
%
% Note: This is automatically set by memoir class. Nevertheless \OnehalfSpacing 
% enables double spacing but leaves single spaced for captions for instance. 
\OnehalfSpacing 
%
% Sets numbering division level
\setsecnumdepth{subsection} 
\maxsecnumdepth{subsubsection}
%
% Chapter style (taken and slightly modified from Lars Madsen Memoir Chapter 
% Styles document
\usepackage{calc,soul,fourier}
\makeatletter 
\newlength\dlf@normtxtw 
\setlength\dlf@normtxtw{\textwidth} 
\newsavebox{\feline@chapter} 
\newcommand\feline@chapter@marker[1][4cm]{%
	\sbox\feline@chapter{% 
		\resizebox{!}{#1}{\fboxsep=1pt%
			\colorbox{gray}{\color{white}\thechapter}% 
		}}%
		\rotatebox{90}{% 
			\resizebox{%
				\heightof{\usebox{\feline@chapter}}+\depthof{\usebox{\feline@chapter}}}% 
			{!}{\scshape\so\@chapapp}}\quad%
		\raisebox{\depthof{\usebox{\feline@chapter}}}{\usebox{\feline@chapter}}%
} 
\newcommand\feline@chm[1][4cm]{%
	\sbox\feline@chapter{\feline@chapter@marker[#1]}% 
	\makebox[0pt][c]{% aka \rlap
		\makebox[1cm][r]{\usebox\feline@chapter}%
	}}
\makechapterstyle{daleifmodif}{
	\renewcommand\chapnamefont{\normalfont\Large\scshape\raggedleft\so} 
	\renewcommand\chaptitlefont{\normalfont\Large\bfseries\scshape} 
	\renewcommand\chapternamenum{} \renewcommand\printchaptername{} 
	\renewcommand\printchapternum{\null\hfill\feline@chm[2.5cm]\par} 
	\renewcommand\afterchapternum{\par\vskip\midchapskip} 
	\renewcommand\printchaptertitle[1]{\color{gray}\chaptitlefont\raggedleft ##1\par}
} 
\makeatother 
\chapterstyle{daleifmodif}
%
% AVV guidelines:
%
% The pages should be numbered consecutively at the bottom centre of the
% page.
\makepagestyle{myvf} 
\makeoddfoot{myvf}{}{\thepage}{} 
\makeevenfoot{myvf}{}{\thepage}{} 
\makeheadrule{myvf}{\textwidth}{\normalrulethickness} 
\makeevenhead{myvf}{\small\textsc{\leftmark}}{}{} 
\makeoddhead{myvf}{}{}{\small\textsc{\rightmark}}
\pagestyle{myvf}
%
% Oscar's command (it works):
% Fills blank pages until next odd-numbered page. Used to emulate single-sided
% frontmatter. This will work for title, abstract and declaration. Though the
% contents sections will each start on an odd-numbered page they will
% spill over onto the even-numbered pages if extending beyond one page
% (hopefully, this is ok).
\newcommand{\clearemptydoublepage}{\newpage{\thispagestyle{empty}\cleardoublepage}}
%
%
% Creates indexes for Table of Contents, List of Figures, List of Tables and Index
\makeindex
% \printglossaries below creates a list of abbreviations. \gls and related
% commands are then used throughout the text, so that latex can automatically
% keep track of which abbreviations have already been defined in the text.
%
% The import command enables each chapter tex file to use relative paths when
% accessing supplementary files. For example, to include
% chapters/brewing/images/figure1.png from chapters/brewing/brewing.tex we can
% use
% \includegraphics{images/figure1}
% instead of
% \includegraphics{chapters/brewing/images/figure1}
\usepackage{import}

% Add other packages needed for chapters here. For example:
\usepackage{lipsum}					%Needed to create dummy text
\usepackage{amsfonts} 					%Calls Amer. Math. Soc. (AMS) fonts
\usepackage[centertags]{amsmath}			%Writes maths centred down
\usepackage{stmaryrd}					%New AMS symbols
\usepackage{amssymb}					%Calls AMS symbols
\usepackage{amsthm}					%Calls AMS theorem environment
\usepackage{newlfont}					%Helpful package for fonts and symbols
\usepackage{layouts}					%Layout diagrams
\usepackage{graphicx}					%Calls figure environment
\usepackage{longtable,rotating}			%Long tab environments including rotation. 
\usepackage[applemac]{inputenc}			%Needed to encode non-english characters 
									%directly for mac
\usepackage{colortbl}					%Makes coloured tables
\usepackage{wasysym}					%More math symbols
\usepackage{mathrsfs}					%Even more math symbols
\usepackage{float}						%Helps to place figures, tables, etc. 
\usepackage{verbatim}					%Permits pre-formated text insertion
\usepackage{upgreek }					%Calls other kind of greek alphabet
\usepackage{latexsym}					%Extra symbols
\usepackage[square,numbers,
		     sort&compress]{natbib}		%Calls bibliography commands 
\usepackage{url}						%Supports url commands
\usepackage{etex}						%eTeXÕs extended support for counters
\usepackage{fixltx2e}					%Eliminates some in felicities of the 
									%original LaTeX kernel
\usepackage[spanish,english]{babel}		%For languages characters and hyphenation
\usepackage{color}                    				%Creates coloured text and background
\usepackage[colorlinks=true,
		     allcolors=black]{hyperref}              %Creates hyperlinks in cross references
\usepackage{memhfixc}					%Must be used on memoir document 
\usepackage{hyperref}									
\usepackage{enumerate}					%For enumeration counter
\usepackage{footnote}					%For footnotes
\usepackage{microtype}					%Makes pdf look better.
\usepackage{rotfloat}					%For rotating and float environments as tables, 
									%figures, etc. 
\usepackage{alltt}						%LaTeX commands are not disabled in 
									%verbatim-like environment
\usepackage[version=0.96]{pgf}			%PGF/TikZ is a tandem of languages for producing vector graphics from a 
\usepackage{tikz}						%geometric/algebraic description.
\usetikzlibrary{arrows,shapes,snakes,
		       automata,backgrounds,
		       petri,topaths}				%To use diverse features from tikz		
%							
%Reduce widows  (the last line of a paragraph at the start of a page) and orphans 
% (the first line of paragraph at the end of a page)
\widowpenalty=1000
\clubpenalty=1000
%
% New command definitions for my thesis
%
\newcommand{\keywords}[1]{\par\noindent{\small{\bf Keywords:} #1}} %Defines keywords small section
\newcommand{\parcial}[2]{\frac{\partial#1}{\partial#2}}                             %Defines a partial operator
\newcommand{\vectorr}[1]{\mathbf{#1}}                                                        %Defines a bold vector
\newcommand{\vecol}[2]{\left(                                                                         %Defines a column vector
	\begin{array}{c} 
		\displaystyle#1 \\
		\displaystyle#2
	\end{array}\right)}
\newcommand{\mados}[4]{\left(                                                                       %Defines a 2x2 matrix
	\begin{array}{cc}
		\displaystyle#1 &\displaystyle #2 \\
		\displaystyle#3 & \displaystyle#4
	\end{array}\right)}
\newcommand{\pgftextcircled}[1]{                                                                    %Defines encircled text
    \setbox0=\hbox{#1}%
    \dimen0\wd0%
    \divide\dimen0 by 2%
    \begin{tikzpicture}[baseline=(a.base)]%
        \useasboundingbox (-\the\dimen0,0pt) rectangle (\the\dimen0,1pt);
        \node[circle,draw,outer sep=0pt,inner sep=0.1ex] (a) {#1};
    \end{tikzpicture}
}
\newcommand{\range}[1]{\textnormal{range }#1}                                             %Defines range operator
\newcommand{\innerp}[2]{\left\langle#1,#2\right\rangle}                                 %Defines inner product
\newcommand{\prom}[1]{\left\langle#1\right\rangle}                                         %Defines average operator
\newcommand{\tra}[1]{\textnormal{tra} \: #1}                                                       %Defines trace operator
\newcommand{\sign}[1]{\textnormal{sign\,}#1}                                                   %Defines sign operator
\newcommand{\sech}[1]{\textnormal{sech} #1}                                                  %Defines sech
\newcommand{\diag}[1]{\textnormal{diag} #1}                                                    %Defines diag operator
\newcommand{\arcsech}[1]{\textnormal{arcsech} #1}                                       %Defines arcsech
\newcommand{\arctanh}[1]{\textnormal{arctanh} #1}                                         %Defines arctanh
%Change tombstone symbol
\newcommand{\blackged}{\hfill$\blacksquare$}
\newcommand{\whiteged}{\hfill$\square$}
\newcounter{proofcount}
\renewenvironment{proof}[1][\proofname.]{\par
 \ifnum \theproofcount>0 \pushQED{\whiteged} \else \pushQED{\blackged} \fi%
 \refstepcounter{proofcount}
 \normalfont 
 \trivlist
 \item[\hskip\labelsep
       \itshape
   {\bf\em #1}]\ignorespaces
}{%
 \addtocounter{proofcount}{-1}
 \popQED\endtrivlist
}
%
%
% New definition of square root:
% it renames \sqrt as \oldsqrt
\let\oldsqrt\sqrt
% it defines the new \sqrt in terms of the old one
\def\sqrt{\mathpalette\DHLhksqrt}
\def\DHLhksqrt#1#2{%
\setbox0=\hbox{$#1\oldsqrt{#2\,}$}\dimen0=\ht0
\advance\dimen0-0.2\ht0
\setbox2=\hbox{\vrule height\ht0 depth -\dimen0}%
{\box0\lower0.4pt\box2}}
%
% My caption style
\newcommand{\mycaption}[2][\@empty]{
	\captionnamefont{\scshape} 
	\changecaptionwidth
	\captionwidth{0.9\linewidth}
	\captiondelim{.\:} 
	\indentcaption{0.75cm}
	\captionstyle[\centering]{}
	\setlength{\belowcaptionskip}{10pt}
	\ifx \@empty#1 \caption{#2}\else \caption[#1]{#2}
}
%
% My subcaption style
\newcommand{\mysubcaption}[2][\@empty]{
	\subcaptionsize{\small}
	\hangsubcaption
	\subcaptionlabelfont{\rmfamily}
	\sidecapstyle{\raggedright}
	\setlength{\belowcaptionskip}{10pt}
	\ifx \@empty#1 \subcaption{#2}\else \subcaption[#1]{#2}
}
%
%An initial of the very first character of the content
\usepackage{lettrine}
\newcommand{\initial}[1]{%
	\lettrine[lines=3,lhang=0.33,nindent=0em]{
		\color{gray}
     		{\textsc{#1}}}{}}
%
% Theorem styles used in my thesis
%
\theoremstyle{plain}
\newtheorem{theo}{Theorem}[chapter]
\theoremstyle{plain}
\newtheorem{prop}{Proposition}[chapter]
\theoremstyle{plain}
\theoremstyle{definition}
\newtheorem{dfn}{Definition}[chapter]
\theoremstyle{plain}
\newtheorem{lema}{Lemma}[chapter]
\theoremstyle{plain}
\newtheorem{cor}{Corollary}[chapter]
\theoremstyle{plain}
\newtheorem{resu}{Result}[chapter]
%
% Hyphenation for some words
%
\hyphenation{res-pec-tively}
\hyphenation{mono-ti-ca-lly}
\hyphenation{hypo-the-sis}
\hyphenation{para-me-ters}
\hyphenation{sol-va-bi-li-ty}
%
%
\begin{document}
% AVV guidlines:
%
% Preliminary pages
% 
% The five preliminary pages must be the Title Page, Abstract, Dedication
% and Acknowledgements, Author's Declaration and Table of Contents.
% These should be single-sided.
% 
% Table of contents, list of tables and illustrative material
% 
% The table of contents must list, with page numbers, all chapters,
 % sections and subsections, the list of references, bibliography, list of
% abbreviations and appendices. The list of tables and illustrations
% should follow the table of contents, listing with page numbers the
% tables, photographs, diagrams, etc., in the order in which they appear
% in the text.
% 
\frontmatter
\pagenumbering{roman}
%
%
% File: Title.tex
% Author: Gayathri Manikutty
% Description: Contains the title page
%
% AVV guidelines:
% 
% At the top of the title page, within the margins, the dissertation should give the title and, if 
% necessary, sub-title and volume number. If the dissertation is in a language other than English, the 
% title must be given in that language and in English. The full name of the author should be in the 
% centre of the page. At the bottom centre should be the words ?A dissertation submitted to the 
% University of Bristol in accordance with the requirements for award of the degree of ? in the 
% Faculty of ...?, with the name of the school and month and year of submission. The word count of 
% the dissertation (text only) should be entered at the bottom right-hand side of the page.
%
%
\begin{titlingpage}
\begin{SingleSpace}
\calccentering{\unitlength} 
\begin{adjustwidth*}{\unitlength}{-\unitlength}
\vspace*{13mm}
\begin{center}
\rule[0.5ex]{\linewidth}{2pt}\vspace*{-\baselineskip}\vspace*{3.2pt}
\rule[0.5ex]{\linewidth}{1pt}\\[\baselineskip]
{\HUGE Thesis Proposal Draft}\\[4mm]

{\HUGE Stealth File Systems in Proactive Android Forensic Rom }\\[4mm]
{\Large by }\\[4mm]
{\Large Sudip Hazra }\\[4mm]
{\Large Under the guidance of }\\[4mm]
{\Large Dr.Prabhakar Mateti }\\[4mm]





\vspace{5mm}
{\large\textsc{December 2015}}
\vspace{12mm}
\end{center}
%\begin{flushright}
%{\small Word count: ten thousand and four}
%\end{flushright}
\end{adjustwidth*}
\end{SingleSpace}
\end{titlingpage}
\clearemptydoublepage
%



\begin{titlingpage}
\begin{SingleSpace}
\vspace{-1cm}
	\begin{center}

	{\HUGE Amrita Vishwa Vidyapeetham}\\[4mm]
	{\HUGE Amritapuri, Kollam}\\[4mm]
	\vspace{11mm}
	\includegraphics[scale=0.9]{logos/AVV_colour}\\
	\vspace{6mm}
			\large\textbf{BONAFIDE CERTIFICATE}
	\end{center}%
	  This is to certify that this mini project report entitled \textbf{Stealth File Systems for Proactive Forensics Support in Custom Android ROM}
	  submitted by \textbf{SUDIP HAZRA (Reg.No : AM.EN.P2CSN15027)} in partial fulfillment of
	  the requirements for the award of the Degree of Master of Technology in CYBER SECURITY SYSTEMS AND NETWORKS
	  is a bonafide record of the work carried out under my guidance and supervision at Amrita School
	  of Engineering.
	  
	  \vspace{6.5mm}
	  \begin{center}
	  This mini project report was evaluated by us on. . . . . . . . . . . . . . .
	  	 
	  \end{center}
	  \vfill
	  
	  
	  % Bottom of the page
	  \begin{flushright}
	  Dr.Prabhaker Mateti\\
	  (Project Guide)\\[1.5cm]
	  Dr.Krishnashree Achuthan\\
	  Professor and Head\\
	  \end{flushright}
	  
	  \begin{flushleft}
	  Date:
	  \end{flushleft}
	  
	   
	  	
	  	  
	
	  
	  
	  
	  
	
	
	
	\end{SingleSpace}
	\end{titlingpage}
\clearemptydoublepage
%
%
% file: dedication.tex
% author: V?ctor Bre?a-Medina
% description: Contains the text for thesis dedication
%

\chapter*{Dedication and acknowledgements}
\begin{SingleSpace}
 At the very outset, I would like to give the first honors to\textbf{ Amma, Mata Amritanandamayi Devi} who gave me the wisdom and knowledge to complete this minor project
under her shelter in this esteemed institution.\\
I express my gratitude to my guide,\textbf{ Prof. Prabhaker Mateti, Associate Professor, Computer Science and Engineering, Wright State University, USA}, for his valuable
suggestion and timely feedback during the course of this minor project.\\
I would like to thank \textbf{Dr Krishnashree Achuthan, Professor and Head of the Center
for Cyber Security}, for giving me useful suggestions and his constant encouragement and
guidance throughout the progress of this minor project.\\
I would like to thank my co-guides\textbf{ Mr. Vipin Pavithran, Assistant Professor, Cyber
Security}, Amrita Vishwa Vidyapeetham for giving me useful suggestions and guidance
throughout the progress of this minor project.\\
I express my special thanks to my colleagues who were with me from the starting of the
minor project, for the interesting discussions and the ideas they shared. Thanks also go
to my friends for sharing so many wonderful moments. They helped me not only enjoy
my life, but also enjoy some of the hardness of this minor project.
\end{SingleSpace}
\clearpage
\clearemptydoublepage
%
%
% File: declaration.tex
% Author: Gayathri Manikutty
% Description: Contains the declaration page
%
% AVV guidelines:
%
% Author's declaration
%
% I declare that the work in this minor thesis was carried out in accordance
% with the requirements of the University's Regulations and Code of Practice
% for Research Degree Programmes and that it has not been submitted for any
% other academic award. Except where indicated by specific reference in the
% text, the work is the candidate's own work. Work done in collaboration with,
% or with the assistance of, others, is indicated as such. Any views expressed
% in the dissertation are those of the author.
%
% SIGNED: .............................................................
% DATE:..........................
%
\chapter*{Author's declaration}
\begin{SingleSpace}
\begin{quote}
\initial{I},\textbf{Sudip Hazra (Registered Number AM.EN.P2CSN15027)} hereby declare that this project report,entitled \textbf {Stealth File Systems for Proactive Forensics Support in Custom Android ROM} is a record of the original work done by me under the guidance of \textbf{ Dr.Prabhaker Mateti},Wright State University  was carried out in accordance with the requirements of the University's Regulations and Code of Practice for Research Degree Programmes and that it has not been submitted for any other academic award. Except where indicated by specific reference in the text, the work is the candidate's own work. Work done in collaboration with,or with the assistance of, others, is indicated as such to the best of my knowledge.

\vspace{2.5cm}
\noindent

\begin{minipage}{.45\linewidth}
\begin{flushleft}
Signature of the Student:\\
Date:

\end{flushleft} 
\end{minipage}
\begin{minipage}{.45\linewidth}
\begin{flushright}
\hspace{-0.75cm}\textsc{Place: Amritapuri.}
\end{flushright}
\end{minipage}
\end{quote}
\begin{center}

CounterSigned
\bigskip
\vspace{1.5cm}

Dr.Krishnashree Achuthan\\
Professor and Head\\
Amrita Center for Cyber Security Systems and Networks
 

\end{center}
\end{SingleSpace}
\clearpage
\clearemptydoublepage
%
\newpage
\thispagestyle{plain}

{\centering\bf ABSTRACT\\}\par\vskip 2cm

\singleSpacing
\noindent
Subramanian, Sripriya.          %% last, first name, upper-lower case mix
M.S.  Department of Computer Science and Engineering,
Wright State University,
2003.                           %% this year
Sniffing the Ethernet with High Quality Tools. %% title

\par\vskip 2cm

\doubleSpacing

[This document is an example collection of files intended to help my
students in using LaTeX as they prepare their theses.  An abstract is
typically about one page.  The following is an example of an abstract
of a thesis.]

This thesis is a study of software quality in the limited context of a
class of network software tools called {\em sniffers}. Sniffers are
network monitoring tools used in the administration of network
security. We analyzed five of the hundreds of the existing sniffers to
determine the causes of poor quality and the methods to eliminate the
problems. We subjected the five selected sniffers to both manual
analysis and analysis by software quality assessment tools. We
classify the 1000+ errors so discovered in these sniffers. Based on
these results, we designed and implemented a new sniffer that is
intended as a model of high quality sniffer. The methods applied to
analyze and enhance the quality of software are studied.

Quality assessment software tools fall into two categories: 1) Static
checkers and 2) Dynamic checkers. Static checkers analyze the source
code. Dynamic checkers stand as guards during run-time. We have
collected and used exhaustively the checker tools that are in the open
source archives. In the course of our use, the quality of the software
quality tools is itself analyzed.

This thesis contributes several case studies of open source software
projects. It also contributes a new distributed sniffer to the open
source. The distributed sniffer can monitor multiple networks and
output desired packet details. We designed our sniffer to include a
MySQL based collector program, packet capture programs and viewer
programs. Our sniffer supports multiple capture and viewer
programs. We applied software engineering methods to eliminate the
quality issues associated with software. We use the results of our
analysis of existing sniffers in avoiding poor quality in our
sniffer. We prepared our sniffer for manual audit using documentation
tools. This new sniffer is an example of {\it literate programming}
that is worthy of study by software engineering students.

\newpage





\clearemptydoublepage
%
\renewcommand{\contentsname} {Table of Contents}
\maxtocdepth{subsection}
\tableofcontents*
\addtocontents{toc}{\par\nobreak \mbox{}\hfill{\bf Page}\par\nobreak}
\clearemptydoublepage
%
\listoftables
\addtocontents{lot}{\par\nobreak\textbf{{\scshape Table} \hfill Page}\par\nobreak}
\clearemptydoublepage
%
\listoffigures
\addtocontents{lof}{\par\nobreak\textbf{{\scshape Figure} \hfill Page}\par\nobreak}
\clearemptydoublepage
%
%
% The bulk of the document is delegated to these chapter files in
% subdirectories.
\mainmatter
%
\import{chapters/chapter01/}{chap01.tex}
\clearemptydoublepage
%
\import{chapters/chapter02/}{chap02.tex}
\clearemptydoublepage
%
\import{chapters/chapter03/}{chap03.tex}
\clearemptydoublepage
%
\import{chapters/chapter04/}{chap04.tex}
\clearemptydoublepage
%
\import{chapters/chapter05/}{chap05.tex}
\clearemptydoublepage
%
\import{chapters/chapter06/}{chap06.tex}
\clearemptydoublepage
%
\import{chapters/chapter07/}{chap07.tex}
\clearemptydoublepage
%
\import{chapters/chapter08/}{chap08.tex}
\clearemptydoublepage
% And the appendix goes here
\appendix
\import{chapters/appendices/}{app0A.tex}
\clearemptydoublepage
%
% Apparently the guidelines don't say anything about citations or
% bibliography styles so I guess we can use anything.
\backmatter

\refstepcounter{chapter}
\bibliography{android,cloud-comp,cloud-storage,forensics,misc}
\bibliographystyle{siam}
\clearemptydoublepage
%
% Add index
%\printindex
%   
\end{document}